\lecture{29}{Спектральный анализ консенсуса}

\sourcevideo
    {https://www.youtube.com/playlist?list=PLOzRYVm0a65fhKDS8cYIC7YCuVGJ4FOJx}
    {Week 8: Lecture 29: Consensus and Average Consensus-2}

Рассматривается линейная динамика
\[
    x^{k+1}=Wx^k,
\]
где \(W\in\mathbb R^{N\times N}\) --- неотрицательная матрица
весов. Спектр \(W\) определяет предельное значение и скорость
консенсуса. Во второй части лекции выводятся два тождества для
нечётных функций относительных состояний, необходимые для анализа
непрерывных нелинейных алгоритмов.

\section{Стохастические матрицы и теорема Перрона--Фробениуса}

\begin{definition}[Строчная стохастичность]
Неотрицательная матрица \(W=(w_{ij})\) называется
строчно-стохастической, если
\[
    \sum_{j=1}^{N}w_{ij}=1,
    \qquad i=1,\ldots,N,
\]
или, эквивалентно,
\[
    W\mathbf1=\mathbf1.
\]
\end{definition}

Следовательно, число \(1\) является собственным значением \(W\), а
\(\mathbf1\) --- соответствующим правым собственным вектором.

\begin{theorem}[Спектральный радиус]
Если \(W\) строчно-стохастична, то каждое её собственное значение
\(\lambda\) удовлетворяет
\[
    |\lambda|\le1,
\]
и потому
\[
    \rho(W)=1.
\]
\end{theorem}

\begin{proof}
Для индуцированной нормы
\[
    \lVert W\rVert_\infty
    =
    \max_i\sum_{j=1}^{N}|w_{ij}|
    =1.
\]
Если \(Wv=\lambda v\) и \(v\ne0\), то
\[
    |\lambda|\lVert v\rVert_\infty
    =
    \lVert Wv\rVert_\infty
    \le
    \lVert W\rVert_\infty\lVert v\rVert_\infty,
\]
откуда \(|\lambda|\le1\). Поскольку \(1\) уже принадлежит спектру,
спектральный радиус равен единице.
\end{proof}

Граф матрицы задаётся соглашением
\[
    w_{ij}>0
    \quad\Longleftrightarrow\quad
    \text{агент }i\text{ получает сообщение от агента }j.
\]
Сильная связность этого графа эквивалентна неприводимости
неотрицательной матрицы. По теореме Перрона--Фробениуса для
неприводимой строчно-стохастической матрицы собственное значение
\(1\) алгебраически простое и существует единственный вектор
\(\pi>0\), удовлетворяющий
\[
    \pi^{\mathsf T}W=\pi^{\mathsf T},
    \qquad
    \mathbf1^{\mathsf T}\pi=1.
\]
Он называется стационарным распределением.

Одной сильной связности недостаточно для сходимости. Например,
\[
    W=
    \begin{pmatrix}
        0&1\\
        1&0
    \end{pmatrix}
\]
имеет собственные значения \(1\) и \(-1\), поэтому состояния
периодически меняются местами.

\begin{definition}[Примитивность]
Неотрицательная матрица \(W\) называется примитивной, если существует
\(m\ge1\), для которого \(W^m>0\) покомпонентно.
\end{definition}

Для неприводимой стохастической матрицы примитивность эквивалентна
апериодичности. В этом случае
\[
    \lambda_1(W)=1,
    \qquad
    |\lambda_i(W)|<1,
    \quad i=2,\ldots,N.
\]
На связном неориентированном графе апериодичность, в частности,
обеспечивается положительными диагональными весами \(w_{ii}>0\).

\section{Взвешенный и средний консенсус}

Пусть матрица \(W\) примитивна и строчно-стохастична. Общий
спектральный результат для простого доминирующего собственного
значения даёт
\[
    W^k\longrightarrow \mathbf1\pi^{\mathsf T}.
\]
Следовательно,
\[
    x^k=W^kx^0
    \longrightarrow
    \mathbf1\pi^{\mathsf T}x^0.
\]
Все агенты сходятся к одному значению
\[
    x_\infty=\pi^{\mathsf T}x^0,
\]
то есть достигается взвешенный консенсус.

\begin{definition}[Двойная стохастичность]
Матрица \(W\) называется дважды стохастической, если
\[
    W\mathbf1=\mathbf1,
    \qquad
    \mathbf1^{\mathsf T}W=\mathbf1^{\mathsf T}.
\]
\end{definition}

Столбцовая стохастичность означает сохранение суммы:
\[
    \mathbf1^{\mathsf T}x^{k+1}
    =
    \mathbf1^{\mathsf T}x^k.
\]
В этом случае нормированный левый собственный вектор равен
\[
    \pi=\frac1N\mathbf1.
\]

\begin{theorem}[Средний консенсус]
Если \(W\) примитивна и дважды стохастична, то
\[
    W^k\longrightarrow
    J:=\frac1N\mathbf1\mathbf1^{\mathsf T}
\]
и
\[
    x^k\longrightarrow Jx^0.
\]
Иными словами,
\[
    x_i^k\longrightarrow
    \frac1N\sum_{j=1}^{N}x_j^0
\]
для каждого агента \(i\).
\end{theorem}

Если \(W=W^{\mathsf T}\) и \(W\mathbf1=\mathbf1\), то матрица
автоматически дважды стохастична. Однако симметричность и связность
без апериодичности недостаточны: приведённая выше матрица перестановки
симметрична и дважды стохастична, но не сходится. Поэтому стандартное
достаточное условие включает неотрицательность, симметричность,
связность графа и положительные диагональные веса.

В симметричном случае существует ортогональное разложение
\[
    W=Q\Lambda Q^{\mathsf T},
\]
причём собственный вектор для \(\lambda_1=1\) равен
\[
    q_1=\frac1{\sqrt N}\mathbf1.
\]
Отсюда
\[
    W^k
    =
    q_1q_1^{\mathsf T}
    +
    \sum_{i=2}^{N}\lambda_i^kq_iq_i^{\mathsf T}
    \longrightarrow
    \frac1N\mathbf1\mathbf1^{\mathsf T}.
\]

\section{Скорость сходимости и топология}

Для симметричной матрицы смешивания положим
\[
    \rho_\perp(W)
    =
    \max_{i=2,\ldots,N}|\lambda_i(W)|.
\]
При среднем консенсусе \(\rho_\perp(W)<1\), и
\[
    \lVert W^k-J\rVert_2
    =
    \rho_\perp(W)^k.
\]
Поэтому
\[
    \lVert x^k-Jx^0\rVert
    \le
    \rho_\perp(W)^k
    \lVert x^0-Jx^0\rVert.
\]
Для относительной точности \(\varepsilon\) достаточно
\[
    k
    \ge
    \frac{\log(1/\varepsilon)}{-\log\rho_\perp(W)}.
\]

Топология определяет одновременно скорость смешивания и цену одного
раунда. Для пути \(P_N\)
\[
    |\mathcal E|=N-1,
    \qquad
    \operatorname{diam}(P_N)=N-1,
\]
поэтому обмен дешёв, но информация распространяется медленно. Для
полного графа \(K_N\)
\[
    |\mathcal E|=\frac{N(N-1)}2,
    \qquad
    \operatorname{diam}(K_N)=1,
\]
поэтому смешивание быстрое, но каждый раунд дорог. У кольца
\[
    |\mathcal E|=N,
    \qquad
    \operatorname{diam}(C_N)
    =
    \left\lfloor\frac N2\right\rfloor,
\]
а степень каждой вершины равна двум.

В экспоненциальных топологиях вершина соединяется с узлами на
циклических расстояниях
\[
    1,2,4,8,\ldots.
\]
И степень, и диаметр таких графов могут иметь порядок \(O(\log N)\).
Они дают промежуточный компромисс между кольцом и полным графом.

Для лапласовского консенсуса скорость связана с алгебраической
связностью \(\lambda_2(L)\). Узкие места уменьшают
\(\lambda_2(L)\), а рёбра между удалёнными частями сети обычно её
увеличивают. Поэтому сравнивать алгоритмы следует по полной
коммуникационной стоимости: числу раундов, умноженному на объём
сообщений за один раунд.

\section{Нечётные взаимодействия и сохранение среднего}

Пусть \(A=(a_{ij})\) --- симметричная матрица весов
неориентированного графа и
\[
    \phi\colon\mathbb R^d\to\mathbb R^d,
    \qquad
    \phi(-z)=-\phi(z).
\]
Примерами служат \(\phi(z)=z\), покоординатный знак и в скалярном
случае
\[
    \phi_\alpha(z)
    =
    |z|^\alpha\operatorname{sign}(z),
    \qquad \alpha>0.
\]

\begin{lemma}[Первая формула симметризации]
Для любых состояний \(x_1,\ldots,x_N\)
\[
    \sum_{i=1}^{N}\sum_{j=1}^{N}
    a_{ij}\phi(x_i-x_j)=0.
\]
\end{lemma}

\begin{proof}
Обозначим левую часть через \(S\). После перестановки индексов
\(i\) и \(j\) получаем
\[
    S
    =
    \sum_{i,j}a_{ji}\phi(x_j-x_i)
    =
    -\sum_{i,j}a_{ij}\phi(x_i-x_j)
    =-S.
\]
Здесь использованы симметрия \(a_{ij}=a_{ji}\) и нечётность
\(\phi\). Следовательно, \(S=0\).
\end{proof}

Рассмотрим непрерывную динамику
\[
    \dot x_i
    =
    -\sum_{j\in\mathcal N_i}
    a_{ij}\phi(x_i-x_j).
\]
По первой формуле
\[
    \frac{\mathrm d}{\mathrm dt}
    \sum_{i=1}^{N}x_i
    =0,
\]
а значит,
\[
    \bar x(t)
    =
    \frac1N\sum_{i=1}^{N}x_i(t)
    =
    \bar x(0).
\]
Если динамика достигает консенсуса, общее значение неизбежно равно
начальному среднему.

\section{Симметризация и функция Ляпунова}

\begin{lemma}[Вторая формула симметризации]
Для произвольных векторов \(e_i\in\mathbb R^d\)
\[
\begin{aligned}
    &\sum_{i,j}
    a_{ij}e_i^{\mathsf T}\phi(x_i-x_j)
    \\
    &\qquad=
    \frac12
    \sum_{i,j}
    a_{ij}(e_i-e_j)^{\mathsf T}\phi(x_i-x_j).
\end{aligned}
\]
\end{lemma}

\begin{proof}
Пусть
\[
    T=\sum_{i,j}a_{ij}e_i^{\mathsf T}\phi(x_i-x_j).
\]
После перестановки индексов и использования симметрии и нечётности
\[
    T
    =
    -\sum_{i,j}a_{ij}e_j^{\mathsf T}\phi(x_i-x_j).
\]
Сложение двух представлений даёт требуемое тождество.
\end{proof}

При \(e_i=x_i\)
\[
\begin{aligned}
    &\sum_{i,j}
    a_{ij}x_i^{\mathsf T}\phi(x_i-x_j)
    \\
    &\qquad=
    \frac12
    \sum_{i,j}
    a_{ij}(x_i-x_j)^{\mathsf T}\phi(x_i-x_j).
\end{aligned}
\]
Для \(\phi(z)=z\) правая часть равна
\[
    \frac12\sum_{i,j}a_{ij}\lVert x_i-x_j\rVert^2,
\]
а для покоординатного знака ---
\[
    \frac12\sum_{i,j}a_{ij}\lVert x_i-x_j\rVert_1.
\]

Положим
\[
    e_i=x_i-\bar x,
    \qquad
    V=\frac12\sum_{i=1}^{N}\lVert e_i\rVert^2.
\]
Поскольку среднее сохраняется, \(\dot e_i=\dot x_i\). Тогда
\[
    \dot V
    =
    -\frac12
    \sum_{i,j}
    a_{ij}(x_i-x_j)^{\mathsf T}\phi(x_i-x_j).
\]
Если
\[
    z^{\mathsf T}\phi(z)>0
    \qquad\text{при }z\ne0,
\]
то \(\dot V\le0\), и функция рассогласования не возрастает.

\section{Негладкие динамики и ограничения анализа}

Для знакового взаимодействия
\[
    \dot x_i
    =
    -\sum_{j\in\mathcal N_i}
    a_{ij}\operatorname{sign}(x_i-x_j)
\]
получаем
\[
    \dot V
    =
    -\frac12
    \sum_{i,j}a_{ij}|x_i-x_j|.
\]
Зависимость от первой степени рассогласований создаёт возможность
конечно-временной сходимости. Однако правая часть разрывна на
множестве \(x_i=x_j\), поэтому строгий анализ проводится в смысле
дифференциальных включений Филиппова.

Для степенного взаимодействия
\[
    \phi_\alpha(z)
    =
    |z|^\alpha\operatorname{sign}(z)
\]
имеем
\[
    \dot V
    =
    -\frac12
    \sum_{i,j}a_{ij}|x_i-x_j|^{1+\alpha}.
\]
Показатель \(0<\alpha<1\) используется в конечно-временных схемах,
а взаимодействие с \(\alpha>1\) усиливает движение вдали от
консенсусного множества. Сочетание
\[
    0<\alpha<1<\beta
\]
лежит в основе фиксированно-временных алгоритмов.

Обе формулы симметризации опираются на условие
\(a_{ij}=a_{ji}\). В ориентированных сетях парное сокращение в общем
случае исчезает; тогда применяются весовая сбалансированность,
специальные левые собственные векторы, push-sum или дополнительные
нормирующие переменные. Первая формула отвечает за сохранение
среднего, а вторая переводит производную функции Ляпунова в сумму
относительных рассогласований на рёбрах.
